\section{Conclusion}
\label{sec:conclusion}

Running the network of an AI weather model in fp16, with its
physical state, normalisation and denormalisation kept in fp32, changes
its forecasts by much less than they differ from the truth: for five
models from three architecture families, the global RMSE changes by at
most 1.4\%, equivalent to at most two hours of lead time at day 10, and
the fp16 forecasts are neither smoother nor less calibrated, nor worse
for the storms examined. At the same time, a forecast is 1.3 to 3.5
times faster and uses 1.4 to 3.7 times less energy.
fp16's limited numeric range, its main risk, left a margin of at least a factor
of 38 in every model; this margin is model-specific and easy to check
with the hooks used here, and should be checked for any model before
fp16 is used. Where in the network the rounding matters most depends on
the architecture rather than on general rules, but since the effect of
the whole fp16 network is already small, no part of the network needs
to be protected. We therefore recommend fp16 inference, with the state
in fp32, as the default for research use of AI weather models on a
single GPU. bf16, although it needs no attention to range, offers no
advantage in speed and loses more skill, and simple 8-bit quantisation
is not a substitute.
